\documentclass{article}
\title{On the Conversion of \LaTeX{} to Word}
\author{A. Researcher}
\date{2026}

\newcommand{\R}{\mathbb{R}}
\newcommand{\norm}[1]{\left\| #1 \right\|}

\begin{document}
\maketitle

\begin{abstract}
We present a converter producing native OMML and live fields.
\end{abstract}

\section{Introduction}\label{sec:intro}
Consider a function $f : \R^n \to \R$ with gradient $\nabla f$.\footnote{The
gradient is the vector of partial derivatives.}
As shown in \eqref{eq:energy}, mass and energy are equivalent.

\begin{equation}\label{eq:energy}
E = mc^2
\end{equation}

\begin{theorem}[Triangle inequality]\label{thm:tri}
For all $x, y$, $\norm{x + y} \leq \norm{x} + \norm{y}$.
\end{theorem}

\begin{proof}
Expand the inner product and apply Cauchy--Schwarz.
\end{proof}

The norm $\norm{x}$ satisfies Theorem \ref{thm:tri}.

\subsection{A multi-line derivation}
\begin{align}
(a+b)^2 &= a^2 + 2ab + b^2 \\
        &= a^2 + b^2 + 2ab
\end{align}

\section{Results}\label{sec:results}
See Section \ref{sec:intro} for context, building on \cite{einstein1905}
and the typesetting of \cite{knuth1984}. A summation:
\[
\sum_{i=1}^{n} i = \frac{n(n+1)}{2}.
\]

\begin{itemize}
\item Editable equations.
\item Native styles.
\item \textbf{Live} cross-references.
\end{itemize}

\begin{table}
\begin{tabular}{lc}
\toprule
Feature & Status \\
\midrule
\multicolumn{2}{c}{Core} \\
Math & OMML \\
\bottomrule
\end{tabular}
\caption{Feature matrix.}\label{tab:features}
\end{table}

\bibliographystyle{plain}
\bibliography{refs}

\end{document}
